% !TEX program = uplatex
\documentclass[uplatex,dvipdfmx,12pt,a4paper]{jsarticle}
\usepackage[margin=1in]{geometry}
\usepackage{amsmath,amssymb,mathtools}
\usepackage[dvipdfmx]{hyperref}

\title{Exhaustive Expansion Self-Reference Generation Rule (EESR):\\
Recursive Sunyata Operator}
\author{}
\date{}

\begin{document}
\maketitle

\begin{abstract}
本稿は、自己記述可能な形式体系に対して「自己参照」と「全展開（全ケース列挙）」を適用すると内部固定点が失われ、外部次元へのメタ射影が必要になる、という生成規則をEESRとして最小限の形式で記述する。
厳密定理ではなく、構造を保持したスケッチとして提示する。
\end{abstract}

\section{Formulation Sketch}

\subsection{Setup}
自己記述可能な形式体系を $S$ とする。

\subsection{(A) Exhaustion operator}
\begin{equation}
  \mathrm{Exh}(S)
  \;\coloneqq\;
  \bigcup_{n \in \mathbb{N}} \{\text{all } 2^n \text{ valuations / cases over } S_n\},
\end{equation}
ここで $S_n$ は $S$ の有限部分体系である。

\subsection{(B) Self-reference operator}
\begin{equation}
  \mathrm{Self}(S) \coloneqq S \cup \{\ulcorner S \urcorner\},
\end{equation}
$\ulcorner S\urcorner$ は $S$ のコード（ゲーデル化・自己記述）を表す。

\subsection{(C) Incompleteness trigger}
\begin{equation}
  \mathrm{Exh}(\mathrm{Self}(S))
  \;\Rightarrow\;
  \neg\exists x\in S\;\text{s.t. } x\text{ is a fixed point for } S.
\end{equation}

\subsection{(D) Meta-projection / dimensional escape}
\begin{equation}
  \Phi(S)
  \;\coloneqq\;
  \min\{S^+ \mid S\subset S^+\;\wedge\; S\text{ is representable in } S^+\}.
\end{equation}

生成規則は次式で与えられる：
\begin{equation}
  \boxed{\;S_{k+1} = \Phi\!\bigl(\mathrm{Exh}(\mathrm{Self}(S_k))\bigr)\;}
\end{equation}

\subsection{One-line summary}
\begin{equation}
  \boxed{\;\mathrm{Self}+\mathrm{Exh}\;\Rightarrow\;\text{No internal fixed point}\;\Rightarrow\;\text{Meta-lift}.\;}
\end{equation}

\end{document}